\documentclass[10pt,a4paper]{amsart}
\usepackage[margin=19mm]{geometry}
\usepackage{amsmath,amssymb,mathtools}
\usepackage[scr=rsfs,cal=euler]{mathalfa}
\usepackage{xfrac}
\usepackage[unicode,hidelinks,bookmarks=false]{hyperref}
\newcommand{\ie}{i.e.\ }
\newcommand{\cf}{cf.\ }
\newcommand{\resp}{respectively\ }
\newcommand{\Subsection}{\subsection}
\providecommand{\nobreakdash}{\nobreak\hbox{-}}
\setlength{\emergencystretch}{2em}
\multlinegap0pt
\usepackage{amssymb}
\usepackage{xcolor}
\newtheorem{prop}{Proposition}
\newtheorem{lemma}{Lemma}
\newtheorem{theorem}{Theorem}
\title{Revisiting the Hales--Jewett Theorem}
\author{Arpita Ghosh}
\date{Source: arXiv:2604.23766v2 (CC BY 4.0)}
\begin{document}
\maketitle

\noindent\emph{Source and licence.} Revisiting the Hales--Jewett Theorem. Version arXiv:2604.23766v2 (CC BY 4.0), distributed by arXiv under the Creative Commons Attribution 4.0 International licence, http://creativecommons.org/licenses/by/4.0/, as recorded in the arXiv licence metadata for this exact version and shown on the abstract page https://arxiv.org/abs/2604.23766v2. Full licence notice: \texttt{licenses/arxiv-2604.23766-CC-BY-4.0.txt}.\par\smallskip

\noindent\textit{Editorial scope.} Counted unit: the article's theorem theorem (source0.tex lines 308--310). The article's complete original proof of it is delivered with it (source0.tex lines 312--408, one \texttt{proof} environment of 97 lines). No mathematics is written, repaired or re-derived: the statement and the proof are the article's own lines, in the article's own notation, and no retained line is substituted or re-flowed.  Retained uncounted as the source-local context the proof needs: lines 152--160; lines 216--228.  Nothing was omitted from any retained span, so the package carries no omission record.  No retained line contains a citation, so the wrapper carries no bibliography and no bibliography entry.  No retained line contains a figure environment, an image-inclusion command or drawing code, so no image is delivered and the package declares no asset dependency.  Editorial formatting only, in addition: an AMS \texttt{amsart} preamble that restates the article's own declarations and macros used by the retained lines, namely the environments \texttt{prop}, \texttt{lemma}, \texttt{theorem} on separate counters, together with the standard packages those lines need.  The retained environments print in this document as Proposition 1, Lemma 1, Theorem 1 in order of appearance, whereas the article prints the same items with the numbering of its own full document.  The complete native source of the article as distributed in the arXiv e-print for this exact version is bundled unchanged as \texttt{sources/199188/source0.tex}.
\begin{prop}\label{pindultra}
Let $\sigma\in \Sigma$, and let $\Psi_\sigma:S^k\to T$ be the associated homomorphism given by $\Psi_\sigma(v_1, \cdots , v_k)= \sigma(v_1\bullet \cdots \bullet v_k)$. Then the continuous extension $\widetilde{\Psi_\sigma}:\beta(S^k)\to \beta T$
satisfies
\[
\widetilde{\Psi_\sigma}\bigl(\mathcal{V}^{\otimes k}\bigr)
   =\bigl(\widetilde{\sigma}(\mathcal{V})\bigr)^{\bullet k}
\]
for every ultrafilter $\mathcal{V}\in \beta S$, where $\widetilde{\sigma}$ is the continuous extension of $\sigma: S \rightarrow T$.
\end{prop}

\begin{lemma}\label{Lemma2}
The following statements are equivalent:

\begin{itemize}
    \item[(a)] Assume that $T$ is a nice subsemigroup of $S,$ and $\Sigma$ is a finite set of retractions from $S$ to $T.$ Moreover assume that $T=B_1 \cup \cdots \cup B_r$ is a coloring of $T$ with finitely many colors. Then there is some $v \in R=S \setminus T$ such that $\{\sigma(v) : \sigma \in \Sigma\}$ is monochromatic. 

   \item [(b)]There exists an ultrafilter $\mathcal{U}$ on $S$ such that
\[
\widetilde{\sigma}(\mathcal{U}) =\widetilde{\tau}(\mathcal{U}) \quad \forall \sigma,\tau \in \Sigma.
\]

\end{itemize}
\end{lemma}

\begin{theorem}[Hales--Jewett Theorem]
Assume that $T$ is a nice subsemigroup of $S,$ and $\Sigma$ is a finite set of retractions from $S$ to $T.$ Moreover assume that $T=B_1 \cup \cdots \cup B_r$ is a coloring of $T$ with finitely many colors. Then there is some $v \in R=S \setminus T$ such that $\{\sigma(v) : \sigma \in \Sigma\}$ is monochromatic.
\end{theorem}

\begin{proof}
We will use induction method on cardinality of $\Sigma.$ For $|\Sigma|=1$, this is trivially true. Let us assume that $|\Sigma|=k$ with $\Sigma= \{\sigma_1, \cdots, \sigma_k\}$. 
Let $\mathcal{U}$ be an ultrafiler of $S$ as in the Lemma \ref{Lemma2} and by induction hypothesis set $\mathcal{V} := \widetilde{\sigma}(\mathcal{U})$ for any $\sigma \in \Sigma \setminus  \{\sigma_k\}$. Now consider $\mathcal{W }= \widetilde{\sigma_k}(\mathcal{U}).$

Define ultrafilters
\[
\mathcal{Z}_{i} := \mathcal{W}^{\bullet i} \bullet \mathcal{V}^{\bullet (r+2-i)}, \quad i=1,\dots,r+1.
\]

Each $\mathcal{Z}_i$ contains one of the colors $B_1,\dots,B_r$. Since there are $r+1$ ultrafilters, there exist $p<q$ and $j$ such that
\[
B_j \in \mathcal{Z}_p \cap \mathcal{Z}_q.
\]

Define a set
\[
\Gamma := \{ v \in T : v^{-1} \bullet B_j \in \mathcal{V}^{\bullet (r+2-q)} \} 
.\]
Here $v^{-1} \bullet B_j= \{w \in T : v\bullet w \in B_j\}.$
Since $B_j \in \mathcal{Z}_p$ and $B_j \in \mathcal{Z}_q$, it follows that
\[
\Gamma \in \mathcal{W}^{\bullet p} \bullet \mathcal{V}^{\bullet (q-p)} \quad \text{and} \quad
\Gamma \in \mathcal{W}^{\bullet q} = \mathcal{W}^{\bullet p} \bullet \mathcal{W}^{\bullet (q-p)}.
\]

Define
\[
\Gamma_1 := \{ u \in T : u^{-1} \bullet \Gamma \in \mathcal{V}^{\bullet (q-p)} \},
\]
\[
\Gamma_2 := \{ u \in T : u^{-1}\bullet \Gamma \in \mathcal{W}^{\bullet (q-p)} \}.
\]
Then $\Gamma_1, \Gamma_2 \in \mathcal{W}^{\bullet p}$.
%Since $\mathcal{W}\in\beta W$, the ultrafilter $\mathcal{W}^{\bullet p}$ is also an ultrafilter on $W$; in particular,
%\[
%W\in\mathcal{W}^{\bullet p}.
%\]
Hence $\Gamma_1\cap\Gamma_2\in\mathcal{W}^{\bullet p},$ so this set is nonempty. Choose
\[
u\in \Gamma_1\cap\Gamma_2.
\]
Therefore, we have $u^{-1} \bullet \Gamma \in \mathcal{V}^{\bullet (q-p)} \cap \mathcal{W}^{\bullet (q-p)}.$

Let $m = q-p$ and consider the tensor product ultrafilter $\mathcal{U}^{\otimes m}$ on $S^m$. For each $\sigma \in \Sigma$, define a map $\Psi_\sigma : S^m \to T$ given by the following formula:
\[
\Psi_\sigma(v_1,\dots,v_m) := \sigma(v_1 \bullet \cdots \bullet v_m).
\]
Then by Proposition \ref{pindultra},

\[
\widetilde{\Psi}_\sigma(\mathcal{U}^{\otimes m})
=
\begin{cases}
\widetilde{\sigma}(\mathcal{U})^{\bullet m}
=\mathcal{V}^{\bullet m},
& \text{if } \sigma\in \Sigma\setminus\{\sigma_k\}.\\ \widetilde{\sigma_k}(\mathcal{U})^{\bullet m} = \mathcal{W}^{\bullet m} & \text{if } \sigma=\sigma_k.
\end{cases}
\]

Hence
\[
\bigcap_{\sigma \in \Sigma} \Psi_\sigma^{-1}(u^{-1} \bullet \Gamma) \in \mathcal{U}^{\otimes m}.
\]
Choose $(v_1,\dots,v_m)$ in the above intersection and we have $\Psi_{\sigma}(v_1, \cdots , v_m) \in u^{-1} \bullet \Gamma.$ Now set $w = v_1 \bullet \cdots \bullet v_m$. Then for every $\sigma \in \Sigma$, $\sigma(w) \in u^{-1} \bullet \Gamma,$
so
\[
u \bullet \sigma(w) \in \Gamma.
\]
%Since $\mathcal{V}\in\beta W$, the ultrafilter $\mathcal{V}^{\bullet(r+2-q)}$ is an ultrafilter on $W$, and therefore
%\[
%W\in\mathcal{V}^{\bullet(r+2-q)}.
%\]
For each $\sigma\in\Sigma$, the inclusion $u \bullet \sigma(w)\in\Gamma$ implies that $(u \bullet \sigma(w))^{-1} \bullet B_j\in\mathcal{V}^{\bullet(r+2-q)}.
$
Since $\Sigma$ is finite, it follows that
\[
\bigcap_{\sigma\in\Sigma}(u \bullet \sigma(w))^{-1} \bullet B_j
\in\mathcal{V}^{\bullet(r+2-q)}.
\]
%Consequently,
%\[
%W\cap\bigcap_{\sigma\in\Sigma}(u\sigma(w))^{-1}B_j
%\in\mathcal{V}^{\bullet(r+2-q)}.
%\]
Hence this set is nonempty, and we may choose
\[
t\in \bigcap_{\sigma\in\Sigma}(u \bullet \sigma(w))^{-1} \bullet B_j.
\]



Now define $v:=u \bullet w \bullet t.$ Since $u,t\in T$ and each $\sigma\in\Sigma$ is a retraction onto $T$, we have 
\[
\sigma(v)=\sigma(u) \bullet \sigma(w) \bullet \sigma(t)=u \bullet \sigma(w) \bullet t\in B_j
\]
for every $\sigma\in\Sigma$. Finally, since $R = S \setminus T$ is a two-sided ideal, we may ensure that $v \in R$.
\end{proof}

\end{document}
